\section{What a confusion matrix identifies}\label{sec:identification}
A confusion matrix can be exactly correct for the target domains and still fail
to determine a user comparison. This is an information problem before it is a
sampling problem. Fix the predictions $q$ and the domain weights $w$ induced
by the estimand in equation~\eqref{eq:estimand}. We now ask what a validation
summary establishes about the fixed target $\theta=w^{\mathsf T}y$, without
specifying a stochastic model for how the classifier's mistakes were generated.

Let $P=\{j:q_j=1\}$ and $N=\{j:q_j=0\}$. Suppose the entire target-population
confusion matrix is known: $f$ predictions in $P$ are false positives and $g$
predictions in $N$ are false negatives. This grants more information than ordinary
held-out validation, whose error rates require a transport assumption.
For a finite collection $B$ of weights, let $S_-(B,m)$ and $S_+(B,m)$ denote
the sums of its $m$ smallest and largest elements, with a zero sum when $m=0$.

\begin{proposition}[Sharp bounds from confusion counts]\label{prop:bounds}
Across binary label vectors compatible with $q,f,g$, the smallest and largest
possible errors $\Delta=\widehat\theta-\theta$ are
\begin{align}
 L&=S_-(w_P,f)-S_+(w_N,g),\notag\\
 U&=S_+(w_P,f)-S_-(w_N,g).
\end{align}
Consequently $[\widehat\theta-U,\widehat\theta-L]$ is the smallest interval
containing every compatible value of $\theta$. Both endpoints are attainable.
\end{proposition}
\begin{proof}
Only false positives and false negatives contribute:
$\Delta=\sum_{j\in\mathrm{FP}}w_j-\sum_{j\in\mathrm{FN}}w_j$.
Select $f$ indices from $P$ and $g$ from $N$. The two selections have disjoint
supports and no other restrictions. Selecting the smallest positive-stratum
weights and largest negative-stratum weights minimizes the error; reversing
these selections maximizes it. Assign the selected indices the corresponding
true labels to attain each endpoint.
\end{proof}
The compatible values form a finite set, so the interval is its convex hull;
not every interior value need be attainable. A reported sign is established by
these counts only if the bounding interval excludes zero. The calculation
requires sorting, rather than a general integer-programming solver.

For example, take $q=(1,1,0,0)$, $w=(1,-1,1,-1)$, and $f=g=1$.
The same confusion matrix permits errors of $-2$, $0$, and $2$.
Prevalence, accuracy, FPR, and FNR are identical across those assignments.
Adding correctly classified domains with zero contrast weight makes accuracy
arbitrarily close to one without resolving the comparison. These are domains
that can affect individual exposure but cancel from the chosen contrast.

\begin{proposition}[When counts suffice]\label{prop:sufficient}
In a stratum of $M$ domains sharing the same predicted label, with exactly
$m$ errors, the error sum is
point identified if and only if $m\in\{0,M\}$ or all its contrast weights are
equal. Its identification width obeys
\begin{equation}\label{eq:dispersion}
 S_+(w,m)-S_-(w,m)
 \leq \min(m,M-m)\{\max_jw_j-\min_jw_j\}.
\end{equation}
Widths add across strata. For several comparisons, counts determine every
comparison if and only if their vector of domain weights is constant within
each stratum with an interior error count.
\end{proposition}
\begin{proof}
Constant weights, no errors, and all errors each determine the sum. If
$0<m<M$ and two weights differ, choose an error set containing one of those
indices but not the other. Exchanging them changes the sum, proving necessity.
The difference between two sums of $m$ terms is at most $m$ times the range.
Taking complements gives the corresponding bound with $M-m$ terms. The
vector statement applies the exchange argument to each coordinate.
\end{proof}
For a stratum mixing predicted labels, the corresponding criterion applies to
the signed contributions $(2q_j-1)w_j$, not to $w_j$ alone.
The result specifies what a finer validation table must preserve. Stratifying
by predicted label alone is sufficient only when the contrast attaches equal
weights within each partially erroneous stratum. Partitioning domains by
similar contrast weights can reduce ambiguity; bins chosen for one comparison
need not work for another. Error rates conditional on user demographics are
not automatically available, because the same domain can enter many groups.

\section{Auditing the target domains}\label{sec:audit}
A target-domain audit can convert these bounds into confidence intervals.
The construction below uses established finite-population counting and audit
ideas. Sorting worst-case error contributions and checking them by random
audits has close precedents in election auditing \citep{stark2008audits}.
Weighted finite-population inference with prediction side information is also
an established problem \citep{shekhar2023audits}. The purpose here is to expose
what confusion-count evidence preserves about reusable predictions, not to
claim a new general audit principle or an optimal interval.

Partition the domains into strata $h=1,\ldots,H$ using information known before
labels are audited. Draw a simple random sample without replacement of a fixed
size $n_h$ from the $M_h$ domains in each stratum. Let $K_h$ be its unknown
number of incorrect predictions and $X_h$ the observed audit-error count.
Conditional on the fixed labels and predictions,
\begin{equation}
 X_h\sim\operatorname{Hypergeometric}(M_h,K_h,n_h).
\end{equation}
Choose nonnegative error budgets $\delta_h$ summing to at most $\delta$.
Retain each feasible integer $K$ for which both
$\Pr_K(X_h\leq x_h)\geq\delta_h/2$ and
$\Pr_K(X_h\geq x_h)\geq\delta_h/2$.
Write the smallest and largest retained integers as $\ell_h,u_h$.
A census stratum reveals $K_h=x_h$ exactly and needs no error budget; an
unaudited stratum contributes no count restriction.

Let $O$ be the audited domains and $U_h$ the unaudited domains in stratum $h$.
Set $c_j=(2q_j-1)w_j$ and
$d_O=\sum_{j\in O}w_j(q_j-y_j)$. Define
\begin{align}
 L_A&=d_O+\sum_h\min_{\ell_h-x_h\leq m\leq u_h-x_h}
                 S_-(c_{U_h},m),\notag\\
 U_A&=d_O+\sum_h\max_{\ell_h-x_h\leq m\leq u_h-x_h}
                 S_+(c_{U_h},m),\label{eq:auditbounds}
\end{align}
where the optimization is over integers. Signed weights make this optimization
necessary: an extreme need not occur at the largest admissible error count.

\begin{proposition}[A simultaneous audit certificate]\label{prop:audit}
Under the fixed stratified sampling design, the interval
$[\widehat\theta-U_A,\widehat\theta-L_A]$ covers the fixed true contrast with
probability at least $1-\delta$. The same event guarantees coverage for every
linear contrast of this one label vector. The endpoints are sharp over the
binary configurations satisfying the count intervals and all audited labels.
\end{proposition}
\begin{proof}
Inverting each pair of hypergeometric tails excludes the true $K_h$ with
probability at most $\delta_h$. A union bound therefore includes all true
counts with probability at least $1-\delta$. On that event, the true errors in
$U_h$ number between $\ell_h-x_h$ and $u_h-x_h$. Sorting and choosing a count
in that range gives the extrema in equation~\eqref{eq:auditbounds}; the choices
are independent restrictions across strata and hence jointly attainable.
The true label vector lies in the same configuration set for every choice of
$w$, which establishes simultaneous coverage.
\end{proof}
The labels and predictions are fixed under this design: systematic or correlated
classifier mistakes do not invalidate the guarantee. Coverage is over repeated
audit samples, not conditional on the realized
sampled identities. It requires neither independent classification errors nor
a large-domain Gaussian approximation. The count confidence region can be
conservative, and sharp projection of that region does not make the resulting
confidence interval shortest. The guarantee concerns one outcome; combining
separate audits of several outcomes requires an additional simultaneous-error
allocation. Optional stopping, adaptive strata, and imperfect audit labels
require further assumptions or other audit procedures.

\paragraph{The information discarded by count projection.}
A large audit need not make these intervals narrow. To see this, let all $N$
predictions equal one, half the weights equal $1/N$ and half equal $-1/N$.
Suppose a fraction $p<1/2$ of labels in each half are wrong. The true contrast
is zero. If $n\to\infty$ but $n/N\to0$, an unstratified audit consistently
learns the overall error fraction. Yet unaudited errors can still be placed
almost entirely on either weight sign: the projected interval width converges
to $2p$. A direct estimate of the weighted audit residual instead has standard
error of order $n^{-1/2}$. Separating the two weight signs into strata removes
this particular information loss. The distinction is between estimating a
weighted residual and bounding it after retaining only counts.

For a declared contrast, an ordinary stratified difference estimator uses
$z_j=w_j(y_j-q_j)$ directly:
\begin{equation}\label{eq:stratified}
 \widetilde\theta=\widehat\theta+
 \sum_h\frac{M_h}{n_h}\sum_{j\in O_h}z_j,\qquad
 \widehat V=\sum_h\frac{M_h^2(1-n_h/M_h)}{n_h}s_{z,h}^2,
\end{equation}
where $s_{z,h}^2$ is the sample variance and a census stratum contributes zero
variance. Under independent stratum draws, the estimate and variance estimator
are design unbiased, with at least two audited domains in each noncensus
stratum. Normal intervals need a separate approximation. A conservative
finite-sample alternative applies Hoeffding's without-replacement inequality
\citep{hoeffding1963} to each bounded residual mean and sums the stratum radii.
More efficient weighted audit procedures are available \citep{shekhar2023audits}.
The count certificate is useful for interpreting validation summaries and
covering many comparisons together, not a general replacement for direct
inference about a chosen comparison.
